\paragraph{Justificación del criterio de polímeros usado en la parte 3.}
Sean $a(P)\ge0$ las normas de las actividades con el margen exponencial ya incluido.
La incompatibilidad incluye dos ocurrencias del mismo polímero. Por el argumento
anclado de F04 se cumple $\sum_{D\sim P}a(D)e^{|D|}\le |P|$ con holgura,
y $\sum_P a(P)e^{|P|}\le CS$, reservando también esta unidad en la exponencial.
Para comprobar la cota del logaritmo, introducir una actividad escalar auxiliar
y desarrollar primero formalmente. La función exponencial de los grafos conexos
enumera todos los grafos en las ocurrencias etiquetadas. Para una colección de $q$
ocurrencias, la suma sobre todos los grafos, con peso $-1$ por arista incompatible,
es el indicador de compatibilidad total por expansión del producto.
Por tanto el logaritmo retiene los grafos conexos y el factor $1/q!$.

En un grafo de incompatibilidad con sus aristas totalmente ordenadas, a cada subgrafo
conexo asígnese su árbol de Kruskal. Los subgrafos que dan un árbol $T$ se obtienen
añadiendo cualquier subconjunto de las aristas posteriores a todas las del camino
que une sus extremos en $T$: al procesarlas sus extremos ya están unidos.
Recíprocamente, cualquier arista extra anterior a la máxima del camino cambiaría
el árbol mínimo. La suma de signos sobre este intervalo es cero o tiene módulo uno.
La suma absoluta después de esta cancelación se domina por la suma sobre árboles.
Para árboles de altura a lo sumo $h$ con raíz $P$ y sin el peso de la raíz
el mayorante etiquetado exponencial se itera desde $F_0(P)=1$ mediante
\[
 F_{h+1}(P)=\exp\Big(\sum_{D\sim P}a(D)F_h(D)\Big).
\]
El factorial de los hijos proviene del orden irrelevante de las ramas; con raíz fijada
la serie lleva $1/(q-1)!$, frente a $1/q!$ en la serie sin raíz. Para $q=1$ el árbol es vacío de aristas.
Por inducción $F_h(P)\le e^{|P|}$. La suma sin raíz queda dominada por
$\sum_P a(P)e^{|P|}$ por designación de una raíz. El margen da convergencia
absoluta en un disco escalar mayor que la unidad: la desigualdad se conserva para
$(1+\delta')a(P)$ con $\delta'>0$ fijo al reducir $\eta_F$. Esto vale con actividades en
el álgebra booleana. La exponencial de la serie conexa coincide allí con la
partición por identidad de coeficientes. El margen restante en talla permite
marcar una o dos ocurrencias y efectuar el centrado de la parte 3.
